\section{Actividad de laboratorio III}

	\subsection{Actividad de simulación: coportamiento del diodo en función de la temperatura}
	\sangria{} Circuito de simulación

	\begin{tikzpicture}
	% Paths, nodes and wires:
	\draw (2, 5.25) -- (3.25, 5.25);
	\draw (2, 3) -- (3.25, 3);
	\draw (-0.25, 5.25) to[american resistor] (1, 5.25);
	\draw (1, 5.25) -- (2, 5.25);
	\draw (-1, 4.404) to[american voltage source] (-1, 3.604);
	\draw (2, 3) -- (-1, 3) -| (-1, 3.6);
	\draw (-0.25, 5.25) -| (-1, 4.404);
	\draw (3.259, 4.4) to[full diode] (3.259, 4);
	\draw (3.259, 4.4) |- (3.109, 5.25);
	\draw (3.259, 4) |- (3.109, 3);
	\end{tikzpicture}

		\subsubsection{Conclusiones del cambio de temperatura}
		\sangria{} Podemos observar como a un aumento de temeperatura $20 \degree C$, $100 \degree C$, $125 \degree C $, aumenta la corriente de conducción para un  mismo voltaje aplicado en polarización directa. Esto se debe a que al incrementarse la temperatura, disminuye la barrera de potencial interna en el diodo (tensión umbral), facilitando el paso de la corriente.La corriente en la región de conducción directa se ve intensificada lo que sugiere una menor resistencia en el diodo, el aumento de la temeperatura también afecta a la corriente de fuga en polarización inversa, ya que al aumentar las temeperatura aumenta la generación de portadores minoritarios.
		\sangria{} Podemos concluir que la conductividad del diodo aumenta, pero sus parámetros electricos se vuelven menos precisos, afectando su funcionamiento en aplicaciones reales.

	\subsection{Actividad de simulación:circuitos recortadores con diodos zener}
		\subsubsection{Pequeña introducción a los diodos como recortadores}
		\sangria Para obtener el máximo valor de salida en el circuito se cumple que en cada semiciclo la tensión de recorte es igual a la tensión del diodo en directa $\approx 0,7V$ más la tensión de ruptura del diodo zener en inversa $V_{z}$, esta relación se va a cumplir tanto para el

		\begin{equation}
			V_{recorte}= V_{z} + 0,7V
		\end{equation}
    \sangria{} \textbf{Circuito A}

\begin{tikzpicture}
	% Paths, nodes and wires:
	\draw (2, 4.966) to[full ZZener diode] (2, 4.716);
	\draw (2, 3.213) to[full ZZener diode] (2, 3.663);
	\draw (2, 4.75) -| (2, 3);
	\draw (2, 5.25) -- (3.25, 5.25);
	\draw (2, 3) -- (3.25, 3);
	\draw (-0.991, 3.634) to[sinusoidal current source] (-0.991, 4.384);
	\draw (-0.25, 5.25) to[american resistor] (1, 5.25);
	\draw (-0.991, 4.384) |- (-0.25, 5.25);
	\draw (1, 5.25) -- (2, 5.25);
	\draw (-0.991, 3.634) |- (2, 3);
	\node[shape=rectangle, draw, line width=1pt, inner sep=0, minimum width=-0.035cm, minimum height=-0.035cm] at (3, 4.75){};
	\draw (2, 5.25) -| (2, 5);
	\end{tikzpicture}

	\sangria{} \textbf{Circuito B}

	\begin{tikzpicture}
	% Paths, nodes and wires:
	\draw (2, 4.5) to[full ZZener diode] (2, 5);
	\draw (2, 4) to[full ZZener diode] (2, 3.25);
	\draw (2, 4.75) -| (2, 3);
	\draw (2, 5.25) -- (3.25, 5.25);
	\draw (2, 3) -- (3.25, 3);
	\draw (-0.991, 3.634) to[sinusoidal current source] (-0.991, 4.384);
	\draw (-0.25, 5.25) to[american resistor] (1, 5.25);
	\draw (-0.991, 4.384) |- (-0.25, 5.25);
	\draw (1, 5.25) -- (2, 5.25);
	\draw (-0.991, 3.634) |- (2, 3);
	\node[shape=rectangle, draw, line width=1pt, inner sep=0, minimum width=-0.035cm, minimum height=-0.035cm] at (3, 4.75){};
	\draw (2, 5.25) -| (2, 5);
	\end{tikzpicture}

	\newpage

	\subsection{Actividad de laboratorio:circuitos recortadores con diodos zener}
	\sangria{} Para calcular la resistencia mínima del circuito, calculamos la tensión eficaz de la fuente, para luego calcular ley de kirchoff de tensión en cada malla, los pasos son los siguientes.

	\par\vspace{-2pt}\par

	\begin{align*}
	&\textbf{Circuito A}\\
	&I_{z1} = 49mA\\
	&I_{z2} = 12,5mA\\
	&\text{Valor RMS solo valido para una sonoidal}\\
	&V_{RMS}=\frac{v_{pp}}{2} \cdot \sqrt{2}\\
	&\text{Resolvamos usando LKT}\\
	&-V_{RMS}+V_{R}-V_{z1}-V_{z2}=0\\
	&V_{R}=V_{RMS}+V_{z1}+V_{z2}\\
	&V_{R}=16,9V+5,1V+20V\\
	&V_{R}=42,070V\\
	&\text{Ahora calculamos el valor de $R$ y su potencia $P_{R}$}\\
	&R=\frac{V_{R}}{I_{z1}}\\
	&R=\frac{42,070V}{49mA}\\
	&R=858,85\Omega \\
	&P_{R}=V_{R} \cdot I_{z1}\\
	&P_{R}=42,070V \cdot 49mA\\
	&P_{R}=2,06W\\
	\\
	&\textbf{Circuito B}\\
	&I_{z1} = 37mA\\
	&I_{z2} = 21mA\\
	\\
	&\text{Valor RMS solo valido para una sonoidal}\\
	&V_{RMS}=\frac{v_{pp}}{2} \cdot \sqrt{2}\\
	\\
	&\text{Resolvamos usando LKT}\\
	&-V_{RMS}+V_{R}-V_{z1}+V_{z2}=0\\
	&V_{R}=V_{RMS}+V_{z1}-V_{z2}\\
	&V_{R}=16,9V+6,8V-12V\\
	&V_{R}=11,7V\\
	\end{align*}
	\begin{align*}
	&\text{Ahora calculamos el valor de $R$ y su potencia $P_{R}$}\\
	&R=\frac{V_{R}}{I_{z1}}\\
	&R=\frac{11,7V}{37mA}\\
	&R=316,21\Omega \\
	&P_{R}=V_{R} \cdot I_{z1}\\
	&P_{R}=11,7V \cdot 37mA\\
	&P_{R}=0,43W\\
	\end{align*}


	\sangria{} Con los valores obtenidos, lo más lógico es llevar ambas resistencias a su valor nominal mas cercano, siendo en ambos casos de $1\Omega$ y $2W$

	\subsubsection{Imágenes Actividad III}
	\imagen[Circuito A]{8cm}{./imagenes/recortadores_A.png}
	\imagen[Salida recortada A]{8cm}{./imagenes/recortadores_A_salida.png}
	\imagen[Circuito B]{8cm}{./imagenes/recortadores_B.png}
	\imagen[Salida recortada B]{8cm}{./imagenes/recortadores_B_salida.png}

	\newpage

		\subsubsection{Conclusiones de los diodos zener como recortadores}
		\sangria{} Los diodos zener como recortadores tienen un correcto funcionamiento tanto en la simulación como en la práctica, en estos dos circuitos la salida del osciloscopio y el simulador coinciden con los límites de tensión de los diodos.

	\subsection{Actividad: análisis hoja de datos}
	\sangria{} La actividad nos dice que debemos hallar el significado y el uso de cada una de estas características de los diodos en la hoja de datos, a continuación las interpretraciones:


	\begin{enumerate}
		\item Características eléctricas(Regimenes Máximos)

		\begin{description}
		\item[$V_{RRM}$:] Repetitive peak reverse voltage, es la máxima tensión pico inversa que el diodo puede sorportar.
		\item[$V_{RSM}$:] Non repetitive peak reverse voltage, es la máxima tensión pico no repetitiva que el diodo puede sorportar, esto quiere decir que el diodo puede sorportar estos niveles de tension durante periodos muy cortos de tiempo.
		\item[$I_{FRM}$:] Forwad RMS current, es la corriente eficaz máxima que el diodo puede conducir en polarización directa de forma continua.
		\item[$I_{FSM}$:] Surge forwad current, es la corriente pico máxima que el diodo puede soportar en pulsos cortos
		\end{description}

		\item Características eléctricas(Regimenes característicos)

		\begin{description}
		\item[$V_{BR}$:] Breakdown voltage. Voltaje en el que ocurre la ruptura eléctrica del dispositivo.
		\item[$I_{R}$:] Reverse current. Corriente que circula en el dispositivo cuando está polarizado en inversa.
		\item[$V_{F}$:] Forward voltage. Voltaje directo en la unión cuando conduce corriente.
		\item[$I_{RM}$:] Peak reverse current. Corriente inversa máxima que puede soportar el dispositivo.
		\item[$I_{F}$:] Forward current. Corriente directa que circula por el dispositivo.
		\end{description}

		\item Características térmicas

		\begin{description}
		\item[$T_{J}$:] Junction temperature. Temperatura máxima permitida en la unión del dispositivo.
		\item[$T_{STG}$:] Storage temperature range. Rango de temperatura de almacenamiento en el que el dispositivo no sufre daños.
		\item[$T_{A}$:] Ambient temperature. Temperatura ambiente durante la operación del dispositivo.
		\item[$R_{\theta JC}$:] Junction-to-case thermal resistance. Resistencia térmica entre la unión y la carcasa del dispositivo, importante para la disipación de calor.
		\end{description}

		\item Características de uso

		\begin{description}
		\item[$V_{RMS}$:] RMS voltage. Valor eficaz de la tensión aplicada al dispositivo.
		\item[$P_{tot}$:] Total power dissipation. Potencia total que el dispositivo puede disipar sin sobrecalentarse.
		\item[$T_{rr}$:] Reverse recovery time. Tiempo de recuperación inversa, es el tiempo que tarda el dispositivo en bloquear la conducción inversa después de cambiar de polarización directa a inversa.
		\end{description}

	\end{enumerate}


\inicioCodigo{}
\par\vspace{10pt}\par
\section{Rúbrica}
\begin{table}[h!]
\centering
\renewcommand{\arraystretch}{1.5}
\begin{tabular}{|>{\raggedright}p{8cm}|c|c|}
\hline
\rowcolor[HTML]{D9D9D9}
\textbf{Tarea} & \textbf{Puntuación} & \textbf{Máximo} \\
\hline
\cellcolor[HTML]{E6E6E6} Presentación del informe & 30 \% & \\
\hline
\cellcolor[HTML]{E6E6E6} Interpretación de los parámetros de las hojas de datos. & 20 \% & \\
\hline
\cellcolor[HTML]{E6E6E6} Identificación de las características de un zener y un rectificador en zona de avalancha (inversa). & 25 \% & \\
\hline
\cellcolor[HTML]{E6E6E6} Diseño y cálculo de circuitos con diodos zener. & 10 \% & \\
\hline
\cellcolor[HTML]{E6E6E6} Defensa de las conclusiones. & 15 \% & \\
\hline
\rowcolor[HTML]{C0C0C0} \textbf{Total} & & \textbf{100 \%} \\
\hline
\end{tabular}
\caption{Tareas y puntuaciones del informe}
\end{table}
\finalCodigo{}
